\lampiransection{HASIL PRA-SURVEI TAHAP 1}
\label{app:prasurvey1}

Lampiran ini memuat hasil pra-survei tahap 1 penelitian \textit{Analisis Strategi Pemasaran Rumah Makan Nasi Gerilya pada Platform GrabFood Menggunakan Matriks SWOT}. Pra-survei dilakukan melalui observasi pada platform GrabFood pada tanggal 04 April 2026 dengan menelaah tampilan toko, rating, jumlah ulasan, kisaran harga, menu terpopuler, promo yang tersedia, serta ringkasan ulasan konsumen. Pra-survei juga dilengkapi dengan observasi singkat terhadap dua kompetitor sejenis, yaitu RM Padang Paripurna dan RM Padang Dendeng Batokok. Tujuan pra-survei ini adalah memperoleh gambaran awal mengenai kekuatan, kelemahan, peluang, dan ancaman Rumah Makan Nasi Gerilya sebagai dasar penyusunan instrumen penelitian utama.

\section{Observasi Platform, Kompetitor, dan Perbandingan Awal}

\begin{table}[h!]
    \centering
    \caption{Identitas Observasi Pra-Survei Tahap 1}
    \label{tab:identitas_prasurvey1}
    \begin{tabularx}{\textwidth}{|l|X|}
        \hline
        \textbf{Aspek}    & \textbf{Keterangan}                                                           \\
        \hline
        Nama usaha        & Rumah Makan Nasi Gerilya                                                      \\
        \hline
        Platform          & GrabFood                                                                      \\
        \hline
        Tanggal observasi & 04 April 2026                                                                 \\
        \hline
        Fokus observasi   & Tampilan toko, rating, ulasan, harga, menu teratas, promo, dan opini konsumen \\
        \hline
    \end{tabularx}
\end{table}

\begin{table}[h!]
    \centering
    \caption{Profil dan Data Utama Rumah Makan Nasi Gerilya}
    \label{tab:data_gerilya}
    \begin{tabularx}{\textwidth}{|l|X|}
        \hline
        \textbf{Kategori Data} & \textbf{Fakta / Temuan}                                                                             \\
        \hline
        Rating                 & 4,7/5,0                                                                                             \\
        \hline
        Total ulasan           & 4.187 ulasan                                                                                        \\
        \hline
        Harga terendah         & Rp11.000 (Nasi Putih)                                                                               \\
        \hline
        Harga tertinggi        & Rp72.000 (Nasi Padang Kari Kambing Gerilya)                                                         \\
        \hline
        Menu teratas           & Nasi Padang Rendang Sapi; Nasi Padang Ayam Goreng; Nasi Padang Ayam Bakar; Nasi Bakar Kikil Tunjang \\
        \hline
        Harga menu teratas     & Rp56.000; Rp45.000; Rp45.000; Rp68.000                                                              \\
        \hline
        Promo Grab             & Diskon Rp8.000 Jajan Hemat; Diskon 25\% maksimal Rp100.000; Diskon 15\% Group Order                 \\
        \hline
        Promo bank             & Raya, BTN, BSI, Mega, Mega Syariah, SBI, OCBC, Danamon, Muamalat                                    \\
        \hline
    \end{tabularx}
\end{table}

\begin{table}[h!]
    \centering
    \caption{Ringkasan Opini Konsumen Rumah Makan Nasi Gerilya}
    \label{tab:opini_gerilya}
    \begin{tabularx}{\textwidth}{|X|X|}
        \hline
        \textbf{Opini Positif}                                                                                                                                                    & \textbf{Opini Negatif}                                        \\
        \hline
        Rasa enak, bumbu dan rempah kuat; porsi tergolong besar; bahan segar; kemasan rapi dan bersih; varian menu beragam; harga dinilai terjangkau; ada indikasi repetisi order &
        Rasa dianggap berubah dibanding dulu; pesanan tidak lengkap atau tidak tersedia; lauk utama dinilai kurang besar atau kurang segar; tekstur nasi kadang keras; proporsi nasi dinilai lebih banyak dibanding lauk; suhu makanan bervariasi \\
        \hline
    \end{tabularx}
\end{table}

Tabel~\ref{tab:identitas_prasurvey1}, Tabel~\ref{tab:data_gerilya}, dan Tabel~\ref{tab:opini_gerilya} menunjukkan bahwa hasil observasi awal memuat data toko, eksposur promosi, serta ringkasan opini konsumen terhadap Rumah Makan Nasi Gerilya pada platform GrabFood.

\begin{table}[h!]
    \centering
    \caption{Ringkasan Data Kompetitor pada Platform GrabFood}
    \label{tab:kompetitor_prasurvey1}
    \begin{tabularx}{\textwidth}{|p{3.5cm}|X|X|}
        \hline
        \textbf{Aspek}     & \textbf{RM Padang Paripurna}                                                                                                       & \textbf{RM Padang Dendeng Batokok}                                                                                                        \\
        \hline
        Rating dan ulasan  & 4,7/5,0; 1.119 ulasan                                                                                                              & 4,7/5,0; 1.313 ulasan                                                                                                                     \\
        \hline
        Harga terendah     & Rp7.000 (Nasi Putih)                                                                                                               & Rp5.000 (Teh Pait Dingin)                                                                                                                 \\
        \hline
        Harga tertinggi    & Rp64.000 (Nasi Bungkus + Ayam Goreng + Rendang)                                                                                    & Rp30.000 (Gulai Tunjang)                                                                                                                  \\
        \hline
        Menu teratas       & Nasi Bungkus Ayam Goreng Paripurna; Nasi Bungkus Ayam Pop; Nasi Bungkus Ayam Gulai; Nasi Bungkus Ayam Rendang                      & Nasi Ayam Rendang; Nasi Dendeng Basah; Nasi Rendang; Nasi Ayam Goreng                                                                     \\
        \hline
        Harga menu teratas & Rp39.500; Rp42.000; Rp38.000; Rp38.000                                                                                             & Rp20.000; Rp29.000; Rp29.000; Rp20.000                                                                                                    \\
        \hline
        Promo Grab         & Diskon Rp8.000 Jajan Hemat; Diskon 15\% Group Order; Diskon 20\% sampai Rp50.000                                                   & Diskon Rp8.000 Jajan Hemat; Diskon 15\% Group; Diskon Rp10.000 min belanja Rp160.000; Diskon 17\% sampai Rp50.000                         \\
        \hline
        Promo bank         & Raya, BTN, BSI, Mega, Mega Syariah, SBI, OCBC, Danamon, Muamalat                                                                   & Raya, BTN, BSI, Mega, Mega Syariah, SBI, OCBC, Danamon, Muamalat                                                                          \\
        \hline
        Ringkasan opini    & Konsumen menilai positif konsistensi bumbu dan porsi; terdapat keluhan terkait rasa yang dinilai standar serta porsi beberapa menu & Konsumen menilai positif rasa dan porsi; terdapat keluhan terkait miskomunikasi pesanan seperti pesanan salah atau kuah dan sambal kurang \\
        \hline
    \end{tabularx}
\end{table}

\begin{table}[h!]
    \centering
    \caption{Matriks Perbandingan Data Faktual Antar Usaha}
    \label{tab:perbandingan_usaha}
    \begin{tabularx}{\textwidth}{|l|r|c|r|c|}
        \hline
        \textbf{Nama Usaha} & \textbf{Ulasan} & \textbf{Rating} & \textbf{Termurah} & \textbf{Rentang Harga} \\
        \hline
        Nasi Gerilya        & 4.187           & 4,7/5,0         & Rp11.000          & Rp45.000 -- Rp68.000   \\
        \hline
        RM Paripurna        & 1.119           & 4,7/5,0         & Rp7.000           & Rp38.000 -- Rp42.000   \\
        \hline
        RM Dendeng Batokok  & 1.313           & 4,7/5,0         & Rp5.000           & Rp20.000 -- Rp29.000   \\
        \hline
    \end{tabularx}
\end{table}

Tabel~\ref{tab:kompetitor_prasurvey1} dan Tabel~\ref{tab:perbandingan_usaha} menunjukkan bahwa kedua kompetitor memiliki rating yang sama dengan Rumah Makan Nasi Gerilya, tetapi berada pada tingkat harga yang lebih rendah. Pada saat yang sama, Rumah Makan Nasi Gerilya memiliki jumlah ulasan tertinggi di antara usaha yang diamati.

\begin{table}[h!]
    \centering
    \caption{Pemetaan SWOT Awal Berdasarkan Pra-Survei Tahap 1}
    \label{tab:swot_prasurvey1}
    \begin{tabularx}{\textwidth}{|l|X|}
        \hline
        \textbf{Kategori}      & \textbf{Indikasi Awal}                                                                                                                      \\
        \hline
        \textit{Strengths}     & Rating 4,7/5,0 dengan jumlah ulasan tinggi; menu unggulan terlihat jelas; promo aktif; opini positif menonjol pada rasa, porsi, dan kemasan \\
        \hline
        \textit{Weaknesses}    & Terdapat opini mengenai penurunan rasa; pesanan tidak lengkap; keluhan terhadap proporsi lauk dan nasi; suhu makanan tidak selalu konsisten \\
        \hline
        \textit{Opportunities} & Dukungan promo platform dan mitra pembayaran digital; reputasi merek yang sudah cukup terlihat pada jumlah ulasan dan repetisi order        \\
        \hline
        \textit{Threats}       & Kompetitor menawarkan produk sejenis dengan harga lebih rendah dan rating yang sama; keluhan berulang dapat memengaruhi persepsi pelanggan  \\
        \hline
    \end{tabularx}
\end{table}

\section{Ringkasan dan Fungsi Pra-Survei Tahap 1}

Berdasarkan hasil observasi, pra-survei tahap 1 menunjukkan bahwa Rumah Makan Nasi Gerilya memiliki interaksi pelanggan yang tinggi, variasi promosi yang aktif, dan posisi harga yang lebih premium dibanding kompetitor yang diamati. Pada saat yang sama, data opini konsumen masih memuat catatan mengenai konsistensi rasa, kelengkapan pesanan, dan kualitas produk yang diterima.

Hasil pra-survei tahap 1 digunakan untuk:
\begin{enumerate}[nosep]
    \item merumuskan kerangka wawancara mendalam dengan pihak manajerial;
    \item memvalidasi poin kuesioner utama survei;
    \item mempertajam analisis kondisi objektif internal dan eksternal;
    \item melengkapi landasan analisis matriks SWOT komprehensif.
\end{enumerate}
